% ============================================================
% Project: THE PHILOSOPHY OF VIABILITY
% Book: Book I — Ontology of Viability
% Book ID: VS.B01
% Part: Part I — Identity and Manifestation
% Chapter: Chapter 07 — Formal and Logical Reconstruction
% File: section_04_definitions_and_conceptual_dependencies.tex
% Section: 04 — Definitions and Conceptual Dependencies
% Purpose: To distinguish definitions, explications, conceptual
% dependencies, logical dependencies, ontological dependence,
% and epistemic justification, while preventing circular
% definitions and illegitimate derivations of ontological claims.
% ============================================================

\section{Definitions and Conceptual Dependencies}
\label{sec:definitions-and-conceptual-dependencies}

\subsection{The Problem of Definition in Ontological Reconstruction}

Formal reconstruction requires a disciplined account of the concepts
through which an ontological position is articulated. Its first task is
not simply to assign a formula to every important expression. It must
determine what kind of conceptual work each formulation performs, what
it presupposes, and what follows from it. A definition may establish
how a term is to be used within a theoretical framework, but it does
not thereby establish that the defined concept corresponds to an
independent feature of reality. A conceptual clarification may reveal
the commitments already contained in an argument, but it does not
automatically demonstrate those commitments. A logical dependency may
show that one proposition follows from another, while leaving entirely
open whether the premises are ontologically adequate.

The central problem is therefore one of controlled dependence. An
ontological theory must employ concepts that are sufficiently
articulated to support reasoning, yet it must not allow the order of
their definitions to determine, without argument, the order or
constitution of being itself. The order in which concepts are
introduced is a feature of theoretical exposition. The order in which
ontological conditions obtain is a different matter. Confusing these
orders can make a theory appear to derive its fundamental commitments
from definitions that merely presuppose them.

This section establishes a framework for distinguishing the principal
forms of definition and dependence relevant to the reconstruction of
identity and manifestation. Its governing claim is:

\begin{quote}
A definition regulates or clarifies conceptual use; a dependency
relation specifies how concepts, propositions, or entities are related.
Neither definitional stipulation nor conceptual ordering, by itself,
establishes an ontological relation.
\end{quote}

The distinction is not a merely terminological precaution. It is
necessary for identifying the actual grounds of a theory, locating
circularity, and preventing a formal reconstruction from attributing
more to its definitions than they can legitimately support.

\subsection{Definition, Explication, and Ontological Characterization}

The word ``definition'' is often used for several different operations.
These operations may overlap in practice, but they should not be
treated as interchangeable.

A \emph{stipulative definition} assigns a specified use to an
expression within a particular inquiry or formal system. It determines
how the expression is to be interpreted for the purposes of that
system. Such a definition can be adopted by convention, provided that
its scope and consequences are made explicit. Its authority is
framework-relative: it establishes a rule of use, not an independent
ontological discovery.

An \emph{explication} seeks to make an initially vague, ambiguous, or
insufficiently articulated concept more precise. It does not merely
replace one word with another. It identifies the role that the concept
is intended to perform, distinguishes its relevant uses, and proposes
conditions under which it can be applied consistently. An explication
may improve a theoretical vocabulary, but its adequacy must be assessed
by criteria appropriate to its purpose, including clarity, coherence,
explanatory usefulness, and fidelity to the phenomenon or problem it
addresses.

An \emph{explicit definition} specifies a concept by stating a
condition or set of conditions under which the defined expression is
to apply. In a formal language, it may take the form of a definitional
equivalence. If a new predicate \(P\) is introduced by reference to an
already available expression \(\varphi\), the definition may be
represented schematically as
\[
  \forall x\bigl(P(x) \leftrightarrow \varphi(x)\bigr).
\]
This schema is legitimate only if the language, variables, and
definitional conditions are sufficiently specified. It does not imply
that every ontological concept can be exhaustively captured by a
predicate of this form.

An \emph{implicit definition} characterizes a term through the role it
occupies within a system of relations, principles, or axioms, rather
than by providing a direct equivalence with previously defined terms.
Implicit characterization can be useful when a concept is identified
through a network of constraints. However, the uniqueness and
interpretability of the concept cannot be presumed merely because a
formal system assigns it a place. Different models may satisfy the
same formal constraints while interpreting the relevant terms
differently.

An \emph{ontological characterization} is a claim about the nature,
status, or conditions of an ontological matter. It may be expressed
through definitions, but it is not reducible to a stipulation. If a
theory maintains that identity is inseparable from being, for example,
the claim is not established simply by defining ``identity'' as
``that which is inseparable from being.'' Such a formulation would
encode the intended commitment in the vocabulary rather than provide
an independent reason for accepting it. The ontological claim must be
stated and assessed as a substantive commitment.

These distinctions yield a first methodological rule: the author must
identify whether a proposed definition introduces a convention,
clarifies a concept, characterizes a theoretical role, or asserts
something about reality. A sentence that performs more than one of
these functions must make their respective contributions
distinguishable.

\subsection{Conceptual Primitives and Definitional Limits}

A formal reconstruction cannot define every term through other terms
without remainder. Any system of definitions must begin somewhere:
with expressions treated as primitive for the purposes of that
system, with a background vocabulary, or with concepts whose use is
clarified without being reduced to a more basic definition.

A \emph{conceptual primitive} is a concept that a given reconstruction
does not define through other concepts within its chosen framework.
Primitiveness is therefore relative to a system of articulation. It
does not mean that the concept is intrinsically simple, absolutely
indefinable, epistemically transparent, or ontologically fundamental.
A concept may be primitive in one formal reconstruction and explicated
through a richer conceptual framework in another.

A \emph{formal primitive} is an expression that is not explicitly
defined within a particular formal language. It may be a primitive
predicate, relation symbol, constant, or other basic expression.
Formal primitiveness concerns the syntax and interpretation of a
formal system. It must not be confused with conceptual primitiveness
in a philosophical theory. A symbol can be formally primitive while
the concept it is intended to represent receives extensive
philosophical clarification outside the formal language.

The distinction matters because an ontological theory may need to
preserve concepts whose significance would be distorted by premature
reduction. Being, identity, and determination, for example, must not
be arranged into a definitional chain that silently makes one of them
nothing more than an expression for another. Their relations require
analysis, but analysis does not entail reduction.

The present reconstruction therefore allows three different
procedures:

\begin{enumerate}
  \item A concept may be treated as primitive within a specified
  formal reconstruction, while its use is clarified through
  explication.
  \item A concept may be explicitly defined for a limited technical
  purpose, provided that the definition does not claim to exhaust
  its ontological significance.
  \item A relation between concepts may be stated as a substantive
  theoretical commitment rather than concealed within a definition.
\end{enumerate}

No concept should be declared absolutely primitive merely because the
current reconstruction has not yet supplied a definition for it.
Likewise, no concept should be declared reducible merely because a
formal equivalence can be written between its symbol and a formula.
The philosophical adequacy of such a reduction remains a separate
question.

\subsection{Conceptual, Logical, Ontological, and Epistemic Dependence}

The expression ``depends on'' is especially hazardous in philosophical
analysis because it can indicate several fundamentally different
relations. A reconstruction that uses this expression without
qualification risks treating a relation of explanation as if it were
a relation of existence, or treating a relation of justification as
if it were a relation of constitution.

\paragraph{Conceptual dependence.}

A concept is conceptually dependent on another concept when its
articulation, intelligibility, or application presupposes the latter
within a specified conceptual framework. For example, a proposed
explication of persistence may invoke identity across a temporal
interval. This indicates a dependency within the conceptual
articulation of persistence. It does not, without further argument,
establish that persistence constitutes identity or that identity is
ontologically produced by persistence.

Conceptual dependence concerns the conditions under which a concept
can be adequately articulated. It may be asymmetric, reciprocal, or
part of a larger network. A reciprocal conceptual relation is not
necessarily viciously circular. Two concepts may illuminate different
aspects of a shared problem. The relevant question is whether the
analysis provides independently intelligible points of entry and
whether the reciprocal clarification adds content rather than merely
repeating the same formulation.

\paragraph{Logical dependence.}

A proposition is logically dependent on another proposition when the
former follows from the latter, either directly or through a specified
set of premises and inference rules. If \(p\) entails \(q\), then
\[
  p \models q.
\]
This relation concerns consequence under a given interpretation or
logical system. It does not establish the truth of \(p\), and it does
not independently establish the ontological adequacy of the
vocabulary in which \(p\) and \(q\) are expressed.

Logical dependence must therefore be distinguished from conceptual
dependence. A concept may be needed to formulate a proposition even
when no proposition involving that concept is logically entailed by
another. Conversely, a proposition may be logically derivable from
premises whose concepts are not conceptually reducible to one another.

\paragraph{Ontological dependence.}

An entity, condition, or feature is ontologically dependent on another
when its being, obtaining, or possibility is grounded in, constituted
by, or otherwise requires the latter in a specified ontological sense.
The precise form of dependence must be stated rather than assumed.
Constitution, grounding, realization, causal dependence, and
existential dependence are not automatically identical relations.

An ontological dependence claim is substantive. It cannot be inferred
solely from the order in which concepts are defined, from the order in
which propositions are introduced, or from the fact that one
expression is used to explain another. If a theory asserts that a
certain manifestation depends ontologically on an existent, it must
identify the intended relation and provide the relevant argument or
foundational commitment.

\paragraph{Epistemic or justificatory dependence.}

A claim is epistemically dependent on another claim when its
justification, evidential support, or rational acceptance relies on
the latter. This relation concerns warrant or justification. It does
not necessarily mirror the ontological order of the matters being
discussed. A claim may be epistemically justified through evidence
that is not ontologically prior to its object; conversely, an
ontologically fundamental condition may be difficult to establish
epistemically.

The four forms of dependence can therefore diverge. A concept may
depend conceptually on another without the corresponding entity
depending ontologically on it. A proposition may follow logically
from a premise without that premise being justified. An ontological
dependence may obtain even when it is not recognized or adequately
represented in a conceptual scheme.

\subsection{The Order of Definition and the Order of Being}

The sequence of definitions in a philosophical work is a sequence of
exposition. It is chosen to make the argument intelligible and to
establish the vocabulary needed for later analysis. It may begin with
familiar concepts, proceed from less complex formulations to more
complex ones, or introduce a concept before fully explaining its
relations. None of these editorial choices establishes an
ontological hierarchy.

This distinction is particularly important for the relation between
being and identity. A text may introduce the concept of being before
the concept of identity because the former provides an accessible
starting point for exposition. It may instead begin with identity
because the inquiry is directed toward the conditions under which
something is identifiable as what it is. Neither order establishes
that being has ontological priority over identity, or that identity
has ontological priority over being.

The reconstruction preserves three distinct commitments:

\begin{description}
  \item[Analytical Distinction.]
  Being and identity can be distinguished in analysis. Their
  conceptual articulation need not collapse into a single expression
  or definition.

  \item[Ontological Inseparability.]
  Being and identity are not treated as independently separable
  ontological items. Their inseparability is a substantive
  ontological commitment, not a consequence of the order in which
  their definitions appear.

  \item[Absence of Ontological Priority.]
  Neither being nor identity is granted ontological priority over
  the other merely on the basis of conceptual, definitional, or
  expository order.
\end{description}

These commitments must remain analytically distinct. The ability to
distinguish two concepts does not by itself establish their
ontological inseparability. Their inseparability does not by itself
establish the absence of ontological priority. Nor does the absence
of priority follow merely from the fact that the concepts are
introduced at the same stage of an argument. Each commitment requires
its own statement and its own appropriate grounds.

The same discipline applies to determination. A theory may examine
how identity is articulated in relation to determination, but it
must not define identity as determination and then present their
relation as a discovery. The relation between them is a matter for
ontological analysis. A definition that makes the relation true by
stipulation would obscure rather than resolve the question.

\subsection{Definitional Circularity and Conceptual Reciprocity}

A circular definition occurs when the expression being defined is
explained through a formulation that presupposes the very concept
whose clarification is sought. The simplest pattern is direct
circularity:
\[
  A \text{ is defined by } B,
  \qquad
  B \text{ is defined by } A.
\]
Such a pattern may provide no independent clarification. More
complex circularity can occur when a chain of definitions eventually
returns to its starting point:
\[
  A \longrightarrow B \longrightarrow C
  \longrightarrow A.
\]

Not every reciprocal relation between concepts is circular in this
defective sense. A theory may legitimately argue that two concepts
are mutually relevant, mutually constraining, or jointly necessary
for an adequate account of a phenomenon. The defect arises when the
proposed definitions supply no independent content, criterion, or
argument by which either concept can be understood or assessed.

A practical test is to ask whether each term can be introduced with
sufficient provisional intelligibility before the reciprocal
relation is asserted. If the answer is negative, the theory must
provide an independent point of entry, such as an example, a
primitive conceptual role, a distinction from neighboring concepts,
or an argument that does not presuppose the conclusion.

A second test is to ask whether the reciprocal relation has
substantive consequences. If the relation merely restates that
\(A\) is related to \(B\) and \(B\) is related to \(A\), without
specifying the kind of relation or the conditions under which it
obtains, it does not yet constitute an explanation.

A third test concerns the status of the conclusion. If a theory
defines identity through inseparability from being and then claims
that identity and being are inseparable, the conclusion is
definitional rather than independently established. If, instead,
the theory articulates being and identity separately, states their
inseparability as an ontological commitment, and supplies reasons
for that commitment, the relation is open to philosophical
assessment.

The present reconstruction therefore distinguishes:

\begin{itemize}
  \item \textbf{Vicious definitional circularity:} a definition
  presupposes the concept it is intended to clarify and provides
  no independent route to understanding it.
  \item \textbf{Explanatory reciprocity:} two concepts clarify or
  constrain one another while retaining distinguishable roles.
  \item \textbf{Substantive ontological inseparability:} a claim
  that two aspects cannot be separated in being, which requires
  an ontological argument or an explicitly adopted commitment.
\end{itemize}

The distinction prevents the theory from treating every reciprocal
relation as a defect, while also preventing it from using
``mutual implication'' as a substitute for explanation.

\subsection{Definitions and the Ontological Vocabulary of Part I}

The concepts developed in Part I require different forms of
articulation. Some can receive relatively precise technical
definitions for particular analytical tasks. Others require
explication, relational characterization, or the explicit statement
of theoretical commitments. The following discussion identifies
methodological requirements without claiming to settle every
ontological question in advance.

\paragraph{Being and existent.}

The terms ``being'' and ``existent'' must be used with a declared
scope. A formal reconstruction may introduce a domain of discourse
whose members are treated as existents for the purposes of
quantification. This formal choice does not settle the full
ontological meaning of being, nor does it establish that the domain
exhausts everything that may be ontologically relevant. The relation
between being and existent therefore requires both conceptual
clarification and a careful account of the formal role assigned to
each expression.

\paragraph{Identity.}

Identity must be articulated without treating it as an additional
property attached to an otherwise complete existent. Its role in
the theory concerns the ontological question of what it is for an
existent to be itself in the relevant sense. This formulation is
an initial characterization, not a complete definition. It does
not reduce identity to numerical identity, individuation,
singularity, invariance, persistence, continuity, determination,
or any other neighboring concept.

A formal language may employ an equality symbol, an identity
predicate, or a relation of equivalence for particular purposes.
Those formal devices must not be assumed to exhaust ontological
identity. The distinction between the formal treatment of identity
and the ontological question remains necessary throughout the
reconstruction.

\paragraph{Determination.}

Determination concerns the articulation of what is determinate
about an existent or an ontological situation. Its relation to
identity must be examined rather than built into the definition of
identity. A determination may contribute to an account of how an
existent is specified, but this does not establish that identity
is identical with the totality of its determinations or reducible
to any one of them.

\paragraph{Distinction and difference.}

Distinction concerns the analytical or ontological recognition of
non-identity, differentiation, or separability under a specified
respect. Difference may be used in several senses and must be
qualified where ambiguity matters. Neither term should be used as
an automatic synonym for identity. The fact that two descriptions
differ does not, by itself, establish two ontologically distinct
identities; nor does the plurality of features or relations
automatically establish a plurality of existents.

\paragraph{Relation.}

A relation is to be characterized by its relata, its formal or
ontological type, and the conditions under which it obtains. A
formal relation symbol may be assigned an arity and an interpretation
in a model. This formal specification does not settle whether the
relation corresponds to an ontological relation, a conceptual
connection, or a relation introduced for a limited analytical
purpose.

\paragraph{Manifestation and expression.}

Manifestation and expression must be articulated in a way that
preserves their distinction. A formal definition may specify a
relation between an existent and a manifestation, or between an
expression and what it expresses, but the ontological significance
of these relations cannot be inferred from the formal notation
alone. Their conceptual and ontological roles must be separately
identified.

\paragraph{Unity, plurality, and singularity.}

These concepts require independent articulation. Unity must not
be defined as singularity, aggregate, or composition. The presence
of multiple components, manifestations, or relations does not
by itself establish a plurality of identities. Conversely, a
formal representation of one domain or one object does not settle
the ontological question of unity. The relevant respect and level
of analysis must be made explicit.

\paragraph{Invariance, persistence, and continuity.}

These concepts may be related, but they must not be collapsed
into one another. Any formal characterization of invariance must
identify the transformation, domain, respect, type, and limits
relative to which invariance is claimed. Persistence and continuity
require their own conditions of application. A definition of one
must not silently import the others as if their equivalence had
already been established.

This catalogue is not an exhaustive dictionary. It identifies
where definitional discipline is especially important and indicates
why a single universal pattern of definition would be inadequate
for the conceptual work of Part I.

\subsection{Explicit Definitions, Implicit Characterization, and Formal Translation}

Formal reconstruction often translates philosophical formulations
into explicit definitions, predicates, relations, or axioms. The
translation must preserve the logical role of the original claim.
A sentence that expresses a substantive ontological commitment
should not be converted into a definition merely because a
definitional form is easier to formalize.

Consider the difference between the following schematic forms:
\[
  \text{(D)} \quad P(x) \leftrightarrow \varphi(x)
\]
and
\[
  \text{(A)} \quad \forall x\,\bigl(P(x) \rightarrow Q(x)\bigr).
\]
The first may introduce or explicate the predicate \(P\) through
\(\varphi\). The second asserts a condition relating \(P\) and
\(Q\). The second is not a definition unless the theoretical
framework explicitly assigns it that role and accepts the
consequences of doing so. It may instead be an axiom, proposition,
or hypothesis.

A further distinction concerns implicit characterization. A set
of axioms may constrain the interpretation of a term without
uniquely determining it. If multiple models satisfy the same
axioms, the axioms provide a structural characterization but do
not necessarily yield a unique interpretation. The reconstruction
must therefore distinguish between:

\begin{enumerate}
  \item the introduction of a symbol;
  \item the specification of its formal role;
  \item the constraints imposed on its interpretation;
  \item the claim that the interpretation corresponds to an
  ontological matter;
  \item the argument supporting that correspondence.
\end{enumerate}

These stages may be connected, but none should be silently
substituted for another. In particular, the existence of a
satisfying formal model shows that the formal constraints can be
jointly interpreted under the relevant semantics. It does not,
without further premises, establish that the model describes
reality.

\subsection{Dependency Graphs and the Control of Conceptual Order}

A dependency graph can help display the architecture of a theory,
provided that its edges are typed. A single unqualified arrow
cannot adequately represent conceptual dependence, logical
entailment, ontological grounding, and epistemic justification.
If these relations are placed in one undifferentiated graph, the
visual structure may suggest a hierarchy that the argument has
not established.

For the purposes of this section, a dependency graph should
therefore distinguish at least the following edge types:

\begin{description}
  \item[Conceptual articulation:] one concept is required to
  articulate or clarify another within the framework.
  \item[Logical entailment:] one proposition follows from another
  under specified premises and inference rules.
  \item[Ontological dependence:] one ontological matter is claimed
  to depend on another in a specified sense.
  \item[Justificatory support:] one claim is offered as a reason
  or evidential ground for another.
  \item[Definitional introduction:] an expression is introduced
  through an explicit rule of definition.
\end{description}

These edge types must not be treated as interchangeable. In
particular, definitional introduction does not entail ontological
dependence, and conceptual articulation does not entail logical
consequence. A graph may show that the exposition introduces
identity after being, but that sequence must not be represented
as an ontological grounding relation unless an independent argument
supports it.

The diagram associated with this section is intended as a
methodological map rather than a completed ontology. It shows
how a disciplined reconstruction can organize its conceptual
and justificatory relations while leaving substantive ontological
relations open to argument.

% =============================================================================
% Project : THE_PHILOSOPHY_OF_VIABILITY
% Book    : Book I — Ontology of Viability
% Book ID : VS.B01
% Part    : Part I — Identity and Manifestation (VS.B01.P01)
% Chapter : Chapter 07 — Formal and Logical Reconstruction
% Chapter ID : VS.B01.P01.C07
%
% File    : 4_ASSETS/diagrams/VS.B01.P01.C07.S04.SS02.P001.DG01.typed_dependency_structure.tex
%
% Section / Chapter / Asset:
%   Diagram DG01 — Typed Dependency Structure
%
% Structural Identifier : VS.B01.P01.C07.S04.SS02.P001.DG01
% Asset Type : DG
%
% Purpose :
%   To distinguish definitional, conceptual, logical, ontological, and
%   justificatory relations among the elements of a formal reconstruction.
%
% Insertion point :
%   Section 04, subsection "Dependency Graphs and the Control of Conceptual
%   Order".
%
% Representational limit :
%   The relations shown are methodologically distinct and must not be
%   collapsed into a single undifferentiated dependency relation.
%
% Dependencies :
%   TikZ; libraries used: arrows.meta, positioning. All are loaded by the
%   project preamble; no per-file package or library loading.
%
% =============================================================================

\begin{figure}[htbp]
    \centering
    \hyphenpenalty=10000 \exhyphenpenalty=10000\relax  % no hyphenation inside diagram labels

    \begin{tikzpicture}[
        font=\small,
        node distance=1.0cm and 0.7cm,
        % --- nodes ---
        concept/.style={
            draw=blue!55!black, fill=blue!8, rounded corners=3pt,
            align=center, text width=3.1cm, minimum height=1.6cm,
            inner sep=5pt, line width=0.7pt
        },
        % --- five relation types: distinct by line pattern, weight and colour ---
        rel/.style={-{Stealth[length=2.5mm]}},
        definitional/.style={rel, line width=1.8pt, draw=blue!55!black},
        conceptual/.style={rel, line width=0.9pt, draw=black!75},
        logical/.style={rel, line width=0.9pt, dashed, draw=purple!60!black},
        ontological/.style={rel, line width=1.2pt, densely dotted,
                            draw=green!40!black},
        justificatory/.style={rel, line width=0.9pt,
                              dash pattern=on 5pt off 2pt on 1pt off 2pt,
                              draw=red!55!black},
        legendtext/.style={font=\scriptsize, align=center, text width=11cm}
    ]

    % ---------------------------------------------------------------
    % Grid of concepts
    % ---------------------------------------------------------------
    \node[concept] (vocab) {Conceptual vocabulary};

    \node[concept, below=of vocab] (explication) {Conceptual\\explication};
    \node[concept, left=of explication]  (definition) {Definitional\\introduction};
    \node[concept, right=of explication] (primitive) {Primitive conceptual\\or formal roles};

    \node[concept, below=of definition]  (formal) {Formal expressions\\and definitions};
    \node[concept, below=of explication] (claims) {Substantive claims\\and propositions};
    \node[concept, below=of primitive]   (interpretation)
        {Interpretation\\and ontological\\characterization};

    \node[concept, text width=4.6cm, minimum height=1.2cm, below=of claims] (assessment)
        {Argument, assessment,\\and justificatory support};

    % ---------------------------------------------------------------
    % Typed relations
    % ---------------------------------------------------------------
    \draw[definitional] (vocab) -- (definition);
    \draw[conceptual]   (vocab) -- (explication);
    \draw[conceptual]   (vocab) -- (primitive);

    \draw[definitional] (definition)  -- (formal);
    \draw[conceptual]   (explication) -- (claims);
    \draw[conceptual]   (primitive)   -- (interpretation);

    \draw[logical]       (formal) -- (claims);
    \draw[ontological]   (interpretation) -- (claims);
    \draw[justificatory] (claims) -- (assessment);
    \draw[justificatory] (interpretation) -- (assessment);

    % ---------------------------------------------------------------
    % Legend (one line sample per relation type)
    % ---------------------------------------------------------------
    \node[below=0.8cm of assessment, legendtext] (legend) {
        \tikz[baseline=-0.5ex]\draw[definitional] (0,0) -- (0.8,0);~definitional
        \quad
        \tikz[baseline=-0.5ex]\draw[conceptual] (0,0) -- (0.8,0);~conceptual
        \quad
        \tikz[baseline=-0.5ex]\draw[logical] (0,0) -- (0.8,0);~logical\\[3pt]
        \tikz[baseline=-0.5ex]\draw[ontological] (0,0) -- (0.8,0);~ontological
        \quad
        \tikz[baseline=-0.5ex]\draw[justificatory] (0,0) -- (0.8,0);~justificatory
    };

    \node[below=0.4cm of legend, text width=11cm, align=left,
          font=\footnotesize\itshape] {
        \textbf{Interpretive constraint:}
        The arrows indicate distinct kinds of theoretical connection.
        They do not establish that every connection is necessary,
        that every path is complete, or that formal introduction
        determines ontological order. In particular, definitional
        introduction and conceptual articulation do not by themselves
        establish ontological dependence.
    };

    \end{tikzpicture}

    \caption{A typed map of definitional and conceptual dependencies. The
    relations shown are methodologically distinct and must not be
    collapsed into a single undifferentiated dependency relation.}
    \label{fig:typed-dependency-structure}
\end{figure}

% =============================================================================
% END OF DIAGRAM DG01
% =============================================================================


\subsection{A Method for Constructing and Auditing Definitions}

A definition should be assessed not only by whether it is concise
or formally well-formed, but also by whether it performs the
function assigned to it. The following procedure provides a
minimum audit.

\begin{enumerate}
  \item \textbf{Identify the target.}
  Specify the expression or concept to be defined. Determine
  whether the target is a technical term, an ontological concept,
  a formal symbol, or a term inherited from an established
  philosophical vocabulary.

  \item \textbf{State the purpose.}
  Clarify whether the formulation is intended to stipulate a use,
  explicate an ambiguous concept, introduce a formal expression,
  characterize a theoretical role, or assert an ontological claim.

  \item \textbf{Specify the scope.}
  State the domain, level of analysis, and context within which
  the definition is intended to operate. A definition adequate
  for a formal model may be insufficient as an ontological
  characterization.

  \item \textbf{Identify presuppositions.}
  List the concepts and commitments on which the definition
  relies. Determine whether any of them presupposes the target
  concept in a way that would make the definition circular.

  \item \textbf{Test conceptual independence.}
  Ask whether the defining expression can be understood
  sufficiently for the proposed definition to clarify the target.
  If not, provide an independent explication, example, primitive
  role, or argument.

  \item \textbf{Type the dependencies.}
  Distinguish conceptual articulation, logical entailment,
  ontological dependence, justificatory support, and definitional
  introduction. Do not infer one from another.

  \item \textbf{Test consequences and alternatives.}
  Determine what follows from the definition and whether
  alternative formulations or interpretations remain possible.
  A stipulative definition should not be presented as uniquely
  compelled when other coherent stipulations are available.

  \item \textbf{Assign claim status.}
  Identify whether the result is a definition, foundational
  commitment, axiom, proposition, thesis, hypothesis, argument,
  or open problem. Separately identify its epistemic or
  justificatory status.
\end{enumerate}

The audit should be recorded in a form that makes the role of each
important definition visible. The accompanying table provides a
classification framework for that purpose.

% ============================================================
% Project: THE PHILOSOPHY OF VIABILITY
% Book: Book I — Ontology of Viability
% Book ID: VS.B01
% Part: Part I — Identity and Manifestation
% Chapter: Chapter 07 — Formal and Logical Reconstruction
% Section: 04 — Definitions and Conceptual Dependencies
% Asset: TL01 — Definition and Dependency Audit
% File: VS.B01.P01.C07.S04.SS04.P001.TL01.definition_and_dependency_audit.tex
% Purpose: To classify definitional operations and dependency
% relations for systematic conceptual and philosophical review.
% Placement: Section 04, subsection
% "A Method for Constructing and Auditing Definitions".
% ============================================================

\begin{table}[htbp]
\centering
\small
\setlength{\tabcolsep}{5pt}
\renewcommand{\arraystretch}{1.35}
\begin{tabular}{|p{2.7cm}|p{4.0cm}|p{4.1cm}|p{3.0cm}|}
\hline
\textbf{Category} &
\textbf{What it establishes} &
\textbf{What it does not establish by itself} &
\textbf{Audit question} \\
\hline

Stipulative definition &
A rule assigning a specified use to an expression within a
framework. &
The independent reality or ontological truth of the stipulated
concept. &
Is the scope of the stipulation explicit? \\
\hline

Explication &
A clarified and more precise articulation of a concept. &
That the proposed clarification is uniquely correct or
ontologically exhaustive. &
Does it preserve the relevant conceptual problem? \\
\hline

Explicit definition &
A stated condition or equivalence governing an expression. &
The ontological adequacy of the expression or its interpretation. &
Are the defining terms independently intelligible? \\
\hline

Implicit characterization &
A conceptual or formal role specified through a system of
relations or constraints. &
A unique interpretation unless uniqueness is separately shown. &
Can alternative interpretations satisfy the same constraints? \\
\hline

Conceptual dependence &
A relation of articulation or intelligibility between concepts. &
Ontological dependence or logical entailment. &
What conceptual presupposition is involved? \\
\hline

Logical dependence &
A relation of consequence between propositions under specified
premises and rules. &
The truth or justification of the premises. &
Which inference rules and premises are required? \\
\hline

Ontological dependence &
A substantive relation of grounding, constitution, realization,
or other specified ontological dependence. &
A merely conceptual or expository priority. &
What kind of ontological dependence is asserted? \\
\hline

Epistemic or justificatory dependence &
A relation in which the warrant for one claim relies on another. &
An ontological relation between the matters described. &
What evidence or reason supports the dependent claim? \\
\hline

Conceptual primitive &
A concept left undefined within a specified conceptual
reconstruction. &
Absolute simplicity or ontological fundamentality. &
Is its role sufficiently clarified for the present purpose? \\
\hline

Formal primitive &
An expression left undefined within a formal language. &
The philosophical primitiveness of the represented concept. &
What interpretation is assigned to the expression? \\
\hline

\end{tabular}
\caption{Classification of definitions and dependency relations
for the audit of conceptual and formal reconstruction.}
\label{tab:definition-and-dependency-audit}
\end{table}

% ============================================================
% END OF ASSET TL01 — CHAPTER 07, SECTION 04
% ============================================================

\subsection{Definitions Do Not Prove Their Ontological Presuppositions}

A recurring error in formal reconstruction is to treat a definition
as if it were evidence for the reality of what it defines. The
error can take several forms.

First, a theory may define a term by incorporating a disputed
ontological thesis into the definiens. It then appears to derive
the thesis from the definition, although the thesis was already
presupposed. This is definitional encoding, not proof.

Second, a theory may define a formal predicate and then infer that
the predicate identifies a real property. The inference requires
an account of interpretation and ontological correspondence.
Formal definability alone does not establish such correspondence.

Third, a theory may stipulate a relation and then treat the
stipulated relation as if it had been discovered in reality. A
formal relation can be useful for modeling, comparison, or
reasoning, but its ontological status requires independent
justification.

Fourth, a theory may introduce a concept as primitive and then
claim that its primitiveness proves its ontological fundamentality.
This inference is invalid. A primitive expression is basic
relative to a formal or conceptual system. Ontological
fundamentality is a different and stronger claim.

The appropriate rule is that definitions may establish the
consequences of a chosen vocabulary, but ontological conclusions
must be supported by premises and arguments whose status is
made explicit. Where a definition is used to clarify an adopted
commitment, the commitment remains a commitment. Where a
definition is used to introduce a formal term, the interpretation
of that term remains a further task.

\subsection{The Independence of Functional Form and Epistemic Status}

The classification of definitions and dependencies must also
respect the two-axis account of claim status established for
this project. The functional or logical form of a statement and
its epistemic or justificatory status answer different questions.

A definition is a functional form: it specifies or regulates the
use of an expression. Its epistemic status may be that of a
stipulated or defined claim. A proposition is also a functional
form, but its epistemic status may range from conditionally
supported to argued or derived. An axiom may be adopted as a
foundational commitment within a formal system, while its
ontological justification remains a separate matter.

Accordingly, the following distinctions must be preserved:

\begin{itemize}
  \item A definition is not automatically true of reality merely
  because it is correctly formulated.
  \item An axiom is not an ontological truth merely because a
  formal system takes it as a starting point.
  \item An argument is not valid merely because it is labeled
  an argument; its premises and inferential steps must be assessed.
  \item A proposition is not established merely because it can
  be expressed in formal notation.
  \item A hypothesis does not become a foundational commitment
  merely because later definitions use its vocabulary.
\end{itemize}

This separation is essential when a conceptual system is being
formalized. Formal clarity can reveal the structure of a claim,
but it cannot replace the assessment of its grounds. The status
of each claim must therefore be recorded on both axes rather
than compressed into a single label.

\subsection{Examples of Definitional Discipline}

The following schematic examples illustrate the difference
between legitimate clarification and illegitimate derivation.

\paragraph{Example 1: Identity and determination.}

Suppose a reconstruction introduces the expression
``determination'' to discuss the ways in which an existent is
specified. It would be methodologically improper to define
identity as the totality of determinations and then claim that
identity is reducible to determination. The conclusion would
have been inserted into the definition. A more disciplined
procedure is to articulate the two concepts separately, state
the proposed relation between them, and examine whether the
relation is necessary, sufficient, partial, or contingent.

\paragraph{Example 2: Invariance and persistence.}

Suppose invariance is formally characterized as the preservation
of a specified feature under a transformation. This may provide
a precise definition for a particular formal context. It does
not establish that invariance is persistence, or that an
invariant feature is sufficient for the persistence of an
identity. The transformation, domain, respect, type, and limits
must be specified; the further relation to persistence requires
a separate argument.

\paragraph{Example 3: Unity and composition.}

Suppose a composite is defined as an entity consisting of
multiple components. This technical definition does not
establish that unity is identical with composition, nor that
every unity must be composite. A formal account of components
may illuminate one kind of unity while leaving other forms of
unity conceptually and ontologically open.

\paragraph{Example 4: Numerical equality and ontological identity.}

Suppose a formal language uses the equality symbol \(=\) with
the standard semantics of numerical identity within a domain.
The formal interpretation specifies how the symbol functions
in that language. It does not settle every philosophical
question about ontological identity, the individuation of
entities, or the adequacy of the domain. The formal symbol
must not be allowed to determine the entire ontological
analysis by convention alone.

These examples illustrate a general principle: the formal
precision of a definition is compatible with the philosophical
openness of the relation it helps articulate.

\subsection{Conditions of Adequate Definition}

A definition is adequate relative to its purpose when it meets
the conditions appropriate to the role it is intended to perform.
No single criterion is sufficient for every kind of definition.
Nevertheless, the following conditions provide a general
standard for the present reconstruction:

\begin{enumerate}
  \item \textbf{Role clarity:} the function of the definition
  is identified.
  \item \textbf{Scope clarity:} its domain and limits are stated.
  \item \textbf{Non-circularity:} it does not merely presuppose
  the concept it is intended to clarify.
  \item \textbf{Conceptual discrimination:} it distinguishes
  the target from neighboring concepts where confusion would
  affect the argument.
  \item \textbf{Consistency:} it does not conflict with other
  definitions or commitments in the same framework without
  an explicit account of the conflict.
  \item \textbf{Inferential transparency:} consequences drawn
  from the definition are distinguished from additional
  premises.
  \item \textbf{Ontological restraint:} the definition does
  not claim to establish an ontological fact that requires
  independent argument.
  \item \textbf{Revisability:} where the definition is
  provisional or framework-relative, its status is recorded
  so that later analysis can revise it without concealing
  the revision.
\end{enumerate}

These conditions do not imply that every concept must receive
a final definition before the theory can proceed. Some concepts
may remain provisionally explicated while their relations are
investigated. The requirement is not premature closure but
explicit control over what has and has not been established.

\subsection{Open Problems and Limits of the Present Framework}

The framework developed here establishes methodological
conditions for definitions and dependency claims. It does not
settle every substantive question about the concepts of Part I.
Several problems remain open.

First, the precise ontological characterization of identity
requires further examination. The present framework prevents
its reduction to determination, numerical identity, invariance,
persistence, or other neighboring concepts, but it does not
claim that the complete positive account has been exhausted
by these distinctions.

Second, the relation between being and identity remains a
substantive ontological issue. Analytical distinction,
ontological inseparability, and the absence of ontological
priority must each be assessed on their own grounds. Their
co-presence in a theoretical framework does not make them
definitions of one another.

Third, the status of conceptual primitives remains open to
revision. A concept treated as primitive for one stage of
formal reconstruction may later receive a richer explication.
Such revision does not necessarily undermine the theory,
provided that the earlier role and the later development are
clearly distinguished.

Fourth, the correspondence between formal structures and
ontological structures requires independent treatment.
Formal models can clarify consistency, consequence, and
alternative interpretation. Their ontological significance
depends on additional arguments concerning interpretation
and adequacy.

Fifth, the relations among conceptual dependence, ontological
dependence, and justificatory dependence require case-by-case
analysis. A general dependency vocabulary can prevent
conflation, but it cannot replace the specific arguments
needed to establish particular dependencies.

Finally, the question of whether some ontological concepts
are irreducible in principle, rather than merely primitive
within a current reconstruction, remains unresolved. The
present section does not convert a methodological limit into
a metaphysical conclusion.

\subsection{Concluding Principle}

The reconstruction of an ontological vocabulary requires
definitions, but the work of ontology cannot be reduced to
definition. Definitions make conceptual commitments explicit,
organize formal expression, and clarify the conditions under
which terms are used. Dependency analysis reveals how those
terms and the propositions involving them are connected.
Neither operation, however, substitutes for the argument
required to establish an ontological claim.

The governing principle of this section is therefore:

\begin{quote}
Every definition must disclose its conceptual function and
presuppositions; every dependency claim must identify the
kind of dependence asserted; and every ontological conclusion
must be supported independently of any merely stipulative
definition that incorporates it.
\end{quote}

This principle provides the basis for the next stage of formal
reconstruction: the explicit identification of foundational
commitments and axiomatic candidates. Before such commitments
can be formalized, it must be clear which statements are
definitions, which are substantive assumptions, and which
remain open to argument.

% ============================================================
% END OF SECTION 04 OF CHAPTER 07
% ============================================================